% LOCKED — do not rewrite. Author-signed abstract (2026-09-21).
% Typo fixes only if the author asks. No clock, area, or seconds.

ASIC-based Artificial Intelligence (AI) accelerators have drawn intense interest
and funding in recent years, yet their reduction to practice remained elusive.
I argue that the issue relies in an old, abstraction-based development model
tied to general-purpose computing CPU design.


This paper presents \tv{}, an open-source toolkit for designing {RISC-V} accelerators
for {AI}, including both hardware and software for out-of-the box
support for PyTorch models acceleration.
At the center of the system is an in-order, single-issue, four-stage
32-bit {RISC-V} core supporting {RV32IM} integer extension and {Zve32x} vector extension.
The core can be complemented with memory-mapped datapath accelerators to form a System-on-Chip,
such as an $8{\times}8$ integer General Matrix
Multiply engine.

A just-in-time compiler ingests a {PyTorch} model, extracts
its computational graph, and turns that graph into {C} code that the stock
{RISC-V} {GNU} Compiler Collection can build.
The compiler extends cleanly to new datapath accelerators and to the
pre- and post-processing they may need, such as tensor tiling.

I prove the stack with {YOLOv3-Tiny} vision model running on Tiny-Vedas.
The design is implemented through the Open{ROAD}
{ASIC} flow on {ASAP7}
Process Design Kit ({PDK}), and exercised on an AMD/Xilinx Alveo~{U280} FPGA.
